\chapter{Architettura dei Motori di Ricerca e Indicizzazione}
\label{chap:indicizzazione-query-processing}

\FloatBarrier
\section{Architettura di Sistema e Paradigma a Due Stadi}
\label{sec:architettura-due-stadi}

Un motore di ricerca su scala Web deve conciliare requisiti architetturali apparentemente inconciliabili: memorizzare e scansionare collezioni testuali non strutturate di dimensioni sterminate garantendo al contempo una latenza di risposta rigorosamente interattiva ($\le 50\text{--}100\text{ ms}$) per ciascuna interrogazione utente~\cite{tonellotto2018}. Per soddisfare tali vincoli sotto volumi elevatissimi di traffico concorrente (migliaia di query al secondo, o QPS), il flusso computazionale viene rigidamente disaccoppiato tra processi asincroni massivi ad alto throughput (\textbf{fase offline}) e processi sincroni a latenza vincolata (\textbf{fase online}).

\FloatBarrier
\subsection{Architettura Distribuita del Motore di Ricerca}
\label{subsec:architettura-distribuita}
Nessun singolo server fisico possiede una capacit\`a di memoria o di calcolo sufficiente per gestire l'intero Web. Come formalizzato da Tonellotto et al.~\cite{tonellotto2018}, un moderno motore di ricerca adotta un'infrastruttura di calcolo distribuita organizzata secondo uno schema \textbf{document-partitioned} (partizionamento orizzontale della collezione):
\begin{itemize}
    \item L'immensa collezione documentale $\mathcal{D}$ ($N \sim 10^9\text{--}10^{11}$ documenti) viene suddivisa in $S$ frammenti indipendenti denominati \textbf{shards} (o partizioni).
    \item Ciascun server di indice (\textit{index server}) ospita un sottoinsieme di documenti e mantiene il relativo indice invertito locale.
    \item Un coordinatore centrale, denominato \textbf{Query Broker}, riceve l'interrogazione dell'utente, la distribuisce in parallelo a tutti gli index shards, raccoglie i risultati parziali prodotti da ciascun nodo ed esegue la fusione dei candidati (\textit{Results Merging}) per estrarre la graduatoria globale.
\end{itemize}

\FloatBarrier
\subsection{Pipeline Offline: Indicizzazione e Compressione}
\label{subsec:pipeline-offline}
La pipeline offline opera in background e a intervalli periodici sulla collezione documentale attraverso tre macro-fasi:
\begin{enumerate}
    \item \textbf{Text Processing e Parsing:} I documenti raccolti tramite crawler vengono ripuliti dai marcatori HTML/XML, tokenizzati, normalizzati e depurati da caratteri non informativi secondo le pipeline esaminate nel \chapref{chap:nlp-preprocessing}.
    \item \textbf{Generazione del Vocabolario e dell'Indice:} Viene estratto il lessico $\mathcal{V} = \{t_1, \dots, t_{|\mathcal{V}|}\}$ e costruita la struttura associativa che mappa ciascun termine alle relative occorrenze nei documenti.
    \item \textbf{Compressione Numerica Lossless:} Poich\'e l'overhead dimensionale dell'indice grezzo pu\`o eguagliare o superare la dimensione del testo originario, vengono applicati algoritmi di compressione per interi ($d$-gaps, codici a byte variabile, Elias-Fano, PForDelta) volti a massimizzare la residenza dei dati nei livelli alti della gerarchia di memoria (RAM e cache CPU, come analizzato nel \chapref{chap:efficienza-cpu}).
\end{enumerate}

\FloatBarrier
\subsection{Pipeline Online e Paradigma di Retrieval a Due Stadi}
\label{subsec:pipeline-online-two-stage}
L'elaborazione delle interrogazioni online adotta una rigida \textbf{Two-Stage Retrieval Pipeline} (o architettura a imbuto), volta a massimizzare l'efficienza senza sacrificare l'accuratezza semantica:

\begin{figure}[H]
\centering
\begin{tikzpicture}[
    node distance=0.7cm and 1.2cm,
    block/.style={rectangle, rounded corners=3pt, draw=navy!80!black, fill=navy!8, thick,
                  font=\small, minimum width=5.0cm, minimum height=0.8cm, align=center},
    stage1/.style={rectangle, rounded corners=3pt, draw=darkgreen!80!black, fill=darkgreen!8, thick,
                   font=\small, minimum width=5.0cm, minimum height=0.9cm, align=center},
    stage2/.style={rectangle, rounded corners=3pt, draw=crimson!80!black, fill=crimson!8, thick,
                   font=\small, minimum width=5.0cm, minimum height=0.9cm, align=center},
    output/.style={rectangle, rounded corners=3pt, draw=navy!90!black, fill=navy!15, thick,
                   font=\small\bfseries, minimum width=5.0cm, minimum height=0.8cm, align=center},
    data/.style={cylinder, shape border rotate=90, draw=gray!70!black, fill=gray!10, thick,
                 font=\footnotesize, aspect=0.25, minimum width=2.4cm, minimum height=1.1cm, align=center},
    arr/.style={-Stealth, thick, navy!80!black}
]
    % Componenti Online
    \node[block] (query) {Query Utente $q$};
    \node[block, below=of query] (rewrite) {Query Expansion \& Riscrittura\\{\scriptsize Correzione refusi, sinonimi, acronimi}};
    \node[stage1, below=of rewrite] (s1) {\textbf{1$^\circ$ Stadio: Candidate Retrieval}\\{\scriptsize Scoring additivo veloce (BM25, WAND)}\\{\scriptsize Indice Invertito residente}};
    \node[stage2, below=of s1] (s2) {\textbf{2$^\circ$ Stadio: Re-Ranking (LtR)}\\{\scriptsize Modelli complessi (LambdaMART, Transformer)}\\{\scriptsize Centinaia/migliaia di feature dense}};
    \node[output, below=of s2] (serp) {SERP Finale (Top-$k$, $k \approx 10$)};

    % Sorgenti Dati Offline
    \node[data, left=1.8cm of s1] (invidx) {Indice\\Invertito\\{\scriptsize Compresso}};
    \node[data, left=1.8cm of s2] (featdb) {Feature\\Store\\{\scriptsize PageRank, CTR...}};

    % Collegamenti
    \draw[arr] (query) -- (rewrite);
    \draw[arr] (rewrite) -- (s1);
    \draw[arr] (invidx) -- (s1);
    \draw[arr] (s1) -- node[right, font=\scriptsize\itshape] {$\approx 1\,000\text{--}10\,000$ candidati ($\mathcal{C}$)} (s2);
    \draw[arr] (featdb) -- (s2);
    \draw[arr] (s2) -- (serp);
\end{tikzpicture}
\caption{Architettura della Two-Stage Retrieval Pipeline: filtraggio massivo ad alta efficienza nel primo stadio ed estrazione dei candidati, seguito da riordinamento fine supervisionato nel secondo stadio~\cite{tonellotto2018}.}
\label{fig:two-stage-pipeline}
\end{figure}

\begin{enumerate}
    \item \textbf{Primo Stadio (Candidate Retrieval / Top-$k$):} opera sull'intera collezione ($N \sim 10^9\text{--}10^{11}$ documenti). Impiega algoritmi di dynamic pruning altamente ottimizzati (es. WAND o MaxScore, trattati nel \chapref{chap:elaborazione-topk-wand-maxscore}) su funzioni di scoring additive e scalari (Okapi BM25). L'obiettivo primario \`e l'altissimo throughput con tempo di risposta sub-millisecondo per shard, estraendo un insieme compatto di candidati:
    \begin{equation}
    \mathcal{C} \subset \mathcal{D}, \quad |\mathcal{C}| \approx 1\,000 \text{--} 10\,000 \text{ documenti}.
    \end{equation}
    \item \textbf{Secondo Stadio (Re-Ranking / Heavy Scoring):} opera esclusivamente sull'insieme $\mathcal{C}$. Interroga un archivio di metadati (\textit{Document Features Repository}) per recuperare centinaia di segnali eterogenei (qualit\`a del dominio, PageRank, freschezza temporale, click-through rate storico, prossimit\`a spaziale, affinit\`a densa neurale). Tali feature alimentano modelli di Machine Learning supervisionati (\textit{Learning-to-Rank}, approfonditi nel \chapref{chap:learning-to-rank-fondamenti}) per produrre la graduatoria finale Top-$k$ visualizzata all'utente ($k \approx 10\text{--}20$).
\end{enumerate}

\FloatBarrier
\section{Modelli di Scoring: Prodotto Scalare e Connessione con il Dense Retrieval}
\label{sec:scoring-dotproduct-ponte}

L'elaborazione delle interrogazioni richiede una metrica algebrica solida per graduare i documenti in base alla loro aderenza al bisogno informativo. Questa sezione riesamina i limiti del modello booleano e del coefficiente di Jaccard, dimostrando come il prodotto scalare su vettori pesati costituisca la base unificante sia del retrieval sparso classico sia del moderno dense retrieval neurale.

\FloatBarrier
\subsection{Richiamo dei Limiti del Modello Booleano e di Jaccard}
\label{subsec:limiti-booleano-jaccard}
Nei sistemi di prima generazione, il recupero dei documenti poggiava sul modello booleano esatto. Come formalizzato nella \secref{sec:boolean-vs-ranked}, tale approccio soffre della grave anomalia del \textit{Feast or Famine}: una congiunzione stringente (\texttt{AND}) rischia di svuotare l'insieme soluzione, mentre una disgiunzione (\texttt{OR}) restituisce centinaia di migliaia di risultati non graduati. Analogamente, l'applicazione ingenua del coefficiente insiemistico di Jaccard (analizzato nella \secref{sec:jaccard-scoring}) fallisce poich\'e penalizza arbitrariamente i documenti lunghi e non discrimina tra parole comuni e parole altamente informative.

\FloatBarrier
\subsection{Il Prodotto Scalare come Misura Unificante di Rilevanza}
\label{subsec:dot-product-misura-unificante}
Per superare tali limitazioni, lo Spazio Vettoriale (VSM) e i moderni schemi di ponderazione (formalizzati nel \chapref{chap:retrieval-vsm-bm25}) proiettano documenti e query in vettori pesati, trasformando la misura di pertinenza nel calcolo di un \textbf{prodotto scalare}:
\begin{equation}
\text{Score}(q, d) = \mathbf{q} \cdot \mathbf{d} = \sum_{t \in \mathcal{V}} w(t, q) \cdot w(t, d).
\label{eq:dot-product-general}
\end{equation}
Poich\'e per una query standard i termini non menzionati hanno peso $w(t, q) = 0$, la sommatoria~\eqref{eq:dot-product-general} si restringe ai soli termini condivisi:
\begin{equation}
\text{Score}(q, d) = \sum_{t \in q \cap d} w(t, q) \cdot w(t, d).
\label{eq:dot-product-intersection}
\end{equation}

Questa formulazione evidenzia una profonda convergenza concettuale tra l'Information Retrieval lessicale classico e il moderno retrieval neurale:
\begin{itemize}
    \item \textbf{Lexical Sparse Retrieval:} I vettori appartengono allo spazio $\mathbb{R}^{|\mathcal{V}|}$, dove ciascuna dimensione corrisponde a un lessema univoco del vocabolario ($|\mathcal{V}| \sim 10^7$). I vettori sono fortemente sparsi e lo score \`e calcolato tramite pesi statistici (TF-IDF, Okapi BM25).
    \item \textbf{Dense Neural Retrieval:} Modelli basati su Transformer (es. bi-encoder neurali) mappano testi e query in vettori continui densi di dimensionalit\`a ridotta $\mathbf{v}_q, \mathbf{v}_d \in \mathbb{R}^k$ ($k \approx 768 \text{ o } 1\,536$). La rilevanza semantica \`e determinata dal prodotto scalare o similarit\`a coseno:
    \begin{equation}
    \text{Score}_{\text{neural}}(q, d) = \mathbf{v}_q \cdot \mathbf{v}_d = \sum_{i=1}^{k} v_{q,i} \cdot v_{d,i}.
    \end{equation}
\end{itemize}
In entrambi i paradigmi, l'operazione fondamentale a runtime consiste nell'aggregazione efficiente di prodotti scalari su collezioni documentali su larga scala.

\FloatBarrier
\section{Strutture Dati ad Alte Prestazioni: Forward Index vs. Indice Invertito}
\label{sec:forward-vs-inverted}

Per gestire la sparsit\`a estrema dello spazio termini--documenti, l'architettura del motore di ricerca deve abbandonare le matrici dense a favore di strutture di memorizzazione compressa. Questa sezione confronta le due modalit\`a canoniche di organizzazione dei dati (Forward Index e Indice Invertito), motivando la centralit\`a dell'indice invertito per la risoluzione rapida delle interrogazioni online.

\FloatBarrier
\subsection{La Matrice Termini--Documenti e il Problema della Sparsit\`a}
\label{subsec:matrice-termini-documenti-sparsita}
Consideriamo la matrice globale termini--documenti $M \in \mathbb{R}^{|\mathcal{V}| \times |\mathcal{D}|}$. I parametri dimensionali di un motore di ricerca web reale evidenziano l'impossibilit\`a fisica di allocare e scansionare tale struttura in forma densa:
\begin{itemize}
    \item Dimensione del lessico: $|\mathcal{V}| \sim 10^7$ termini unici.
    \item Dimensione della collezione: $|\mathcal{D}| \sim 10^9 \text{--} 10^{11}$ documenti web.
    \item Numero totale di celle: $|\mathcal{V}| \times |\mathcal{D}| \approx 10^{16} \text{--} 10^{18}$ posizioni.
\end{itemize}
Poich\'e un documento medio contiene tra $10^2$ e $10^3$ vocaboli distinti, la densit\`a di elementi non nulli \`e infinitesima:
\begin{equation}
\frac{\text{elementi non nulli}}{|\mathcal{V}| \times |\mathcal{D}|} \ll 0.001\%.
\end{equation}
Oltre il $99.99\%$ delle celle contiene il valore zero. Memorizzare la matrice in formato denso comporterebbe la saturazione immediata della memoria secondaria e una complessit\`a temporale di valutazione query proibitiva, pari a $\BigO{|\mathcal{V}| \cdot |\mathcal{D}|}$, costringendo la CPU a esaminare miliardi di documenti irrilevanti.

\FloatBarrier
\subsection{Rappresentazione tramite Vettori Sparsi: Forward Index vs. Indice Invertito}
\label{subsec:vettori-sparsi-fwd-inv}
Per eliminare gli zeri, la matrice viene memorizzata come un insieme di vettori sparsi composti da coppie $(\text{Identificatore}, \text{Valore})$. L'orientamento di serializzazione determina due strutture complementari, sintetizzate nella Tabella~\ref{tab:fwd-vs-inv}:
\begin{itemize}
    \item \textbf{Forward Index (Column-Wise / Document-Oriented):} memorizza per ciascun documento la lista ordinata dei termini che vi compaiono, ad esempio $d_1 \to [(t_2, w_2), (t_7, w_7)]$.
    \item \textbf{Indice Invertito (Row-Wise / Term-Oriented):} memorizza per ciascun termine la lista ordinata dei documenti in cui ricorre, ad esempio $t_1 \to [(d_2, w_1), (d_5, w_5)]$.
\end{itemize}

\begin{table}[H]
\centering
\small
\begin{tabularx}{\linewidth}{l >{\raggedright\arraybackslash}X >{\raggedright\arraybackslash}X}
\toprule
\textbf{Propriet\`a} & \textbf{Forward Index} & \textbf{Indice Invertito} \\
\midrule
\textbf{Orientamento} & Per colonna (Document-Oriented). & Per riga (Term-Oriented). \\
\addlinespace
\textbf{Struttura logica} & $\text{DocID} \to [(t_a, \text{tf}_a), \dots]$ & $t \to [(\text{DocID}_i, \text{tf}_i), \dots]$ \\
\addlinespace
\textbf{Complessit\`a query} & $\BigO{|\mathcal{D}|}$: richiede ispezione di tutti i documenti. & $\BigO{\sum_{t \in q} |\mathcal{L}(t)|}$: accede solo ai documenti candidati. \\
\addlinespace
\textbf{Ruolo nel motore} & Memorizzazione metadati, feature store, snippet generation. & Cuore computazionale del First-Stage Retrieval (matching rapido). \\
\bottomrule
\end{tabularx}
\caption{Confronto sistematico tra Forward Index e Indice Invertito.}
\label{tab:fwd-vs-inv}
\end{table}

\FloatBarrier
\section{Anatomia, Layout di Memoria e Compressione dell'Indice Invertito}
\label{sec:anatomia-compressione-indice}

Un indice invertito di livello industriale si articola in due componenti fisiche distinte: il \textbf{Dizionario} (o Lessico $\mathcal{V}$) e l'insieme delle \textbf{Posting List}.

\begin{figure}[H]
\centering
\resizebox{0.95\linewidth}{!}{%
\begin{tikzpicture}[
    lex/.style={rectangle, rounded corners=3pt, draw=navy!80!black, fill=navy!6, thick,
                font=\small, minimum width=4.6cm, minimum height=1.15cm, align=left},
    post/.style={rectangle, rounded corners=3pt, draw=darkgreen!80!black, fill=darkgreen!6, thick,
                 font=\small, minimum width=6.8cm, minimum height=1.15cm, align=left},
    arr/.style={-Stealth, thick, navy!85!black}
]
    % Titoli allineati sulla medesima quota verticale
    \node[font=\footnotesize\bfseries, text=navy!90!black, align=center] at (-3.2, 1.2) {DIZIONARIO LESSICALE\\[2pt]\scriptsize (In RAM / B-Tree / Trie)};
    \node[font=\footnotesize\bfseries, text=darkgreen!90!black, align=center] at (3.2, 1.2) {POSTING LISTS (Stream Disaccoppiati)\\[2pt]\scriptsize (Su Disco / File Invertito / Memoria Secondaria)};

    % Riga 1: caesar
    \node[lex] (l1) at (-3.2, 0.0) {\textbf{``caesar''} \hfill $df=3$ \hfill $\bullet$};
    \node[post] (p1) at (3.2, 0.0) {
        \textbf{DocIDs:} $[1, 2, 4]$ \quad $\to \Delta = [1, 1, 2]$\\
        \textbf{TFs:} \phantom{Doc} $[232, 227, 2]$
    };

    % Riga 2: mercy
    \node[lex] (l2) at (-3.2, -1.6) {\textbf{``mercy''} \hfill $df=3$ \hfill $\bullet$};
    \node[post] (p2) at (3.2, -1.6) {
        \textbf{DocIDs:} $[1, 3, 4]$ \quad $\to \Delta = [1, 2, 1]$\\
        \textbf{TFs:} \phantom{Doc} $[2, 3, 5]$
    };

    % Riga 3: worser
    \node[lex] (l3) at (-3.2, -3.2) {\textbf{``worser''} \hfill $df=2$ \hfill $\bullet$};
    \node[post] (p3) at (3.2, -3.2) {
        \textbf{DocIDs:} $[1, 3]$ \quad\quad $\to \Delta = [1, 2]$\\
        \textbf{TFs:} \phantom{Doc} $[2, 1]$
    };

    % Collegamenti rigorosamente orizzontali
    \draw[arr] (l1.east) -- (p1.west);
    \draw[arr] (l2.east) -- (p2.west);
    \draw[arr] (l3.east) -- (p3.west);
\end{tikzpicture}%
}
\caption{Anatomia di un indice invertito con disaccoppiamento fisico tra stream dei $\text{DocID}$ (compressi tramite $d\text{-gaps}$) e stream delle frequenze $\text{tf}$.}
\label{fig:inverted-index-anatomy}
\end{figure}

\FloatBarrier
\subsection{Il Dizionario (Lexicon)}
\label{subsec:dizionario-lexicon}
Il dizionario memorizza il vocabolario $\mathcal{V}$. Per ciascun termine $t$, registra la stringa lessicale, la Document Frequency $df_t$ (utilizzata per ricavare l'IDF precalcolato), la Collection Frequency $F_t$ e il puntatore di memoria o l'offset su storage secondario che localizza la corrispondente posting list.
Dato che $|\mathcal{V}|$ \`e dell'ordine di milioni di lessemi, il dizionario viene mantenuto interamente residente in \textbf{RAM}, indicizzato tramite strutture di accesso rapido come B-Tree, Trie compressi o tabelle hash perfette.

\FloatBarrier
\subsection{Le Posting List: Disaccoppiamento Fisico degli Stream}
\label{subsec:disaccoppiamento-stream}
Concettualmente, la posting list associata al termine $t$ \`e una sequenza di tuple:
\begin{equation}
\mathcal{L}(t) = \big[ \langle d_1, \text{tf}_{t,d_1} \rangle, \langle d_2, \text{tf}_{t,d_2} \rangle, \dots \big].
\end{equation}
A livello implementativo, memorizzare array continui di strutture composte (es. struct C++ contenenti \texttt{uint32\_t doc\_id} e \texttt{uint32\_t tf}) \`e altamente sub-ottimale. Per massimizzare i fattori di compressione numerica, i motori di ricerca operano il \textbf{disaccoppiamento fisico in stream paralleli}:

\begin{enumerate}
    \item \textbf{Stream dei DocID (Monotono Crescente):}
    \begin{equation}
    \mathbf{D}_t = [\text{DocID}_1, \text{DocID}_2, \dots, \text{DocID}_m] \quad \text{con } \text{DocID}_i < \text{DocID}_{i+1}.
    \end{equation}
    Sfruttando la stretta monotonicit\`a, i valori assoluti vengono convertiti nelle differenze tra identificatori contigui, dette \textbf{$d\text{-gaps}$} o delta:
    \begin{equation}
    \Delta_1 = \text{DocID}_1, \quad \Delta_i = \text{DocID}_i - \text{DocID}_{i-1} \quad (\forall i \ge 2).
    \end{equation}
    I delta generati sono interi piccoli (spesso $\Delta \le 5$ per termini ad alta frequenza), comprimibili ad altissima densit\`a tramite algoritmi dedicati (\textit{Variable Byte}, \textit{Elias-Fano}, \textit{PForDelta}, \textit{Binary Interpolative Coding}~\cite{moffat2000}).
    
    \item \textbf{Stream delle Frequenze (Non Ordinato):}
    \begin{equation}
    \mathbf{F}_t = [\text{tf}(t, d_1), \text{tf}(t, d_2), \dots, \text{tf}(t, d_m)].
    \end{equation}
    Le frequenze locali non godono di ordinamento; tuttavia, per la legge di Zipf, la loro distribuzione empirica \`e fortemente concentrata su valori minimi ($1, 2, 3$). Esse vengono compresse simmetricamente tramite codici a lunghezza variabile senza richiedere la differenziazione delta.
\end{enumerate}

\FloatBarrier
\subsection{Trade-Off Architetturale: Interi Compressi vs. Float Precalcolati}
\label{subsec:trade-off-interi-float}
Una decisione progettuale critica consiste nello stabilire se memorizzare nell'indice il punteggio finale precalcolato $w(t, d) \in \mathbb{R}$ a virgola mobile oppure l'intero grezzo $\text{tf}(t, d) \in \mathbb{N}$:

\begin{table}[H]
\centering
\small
\begin{tabularx}{\linewidth}{l X X}
\toprule
\textbf{Parametro} & \textbf{Punteggio Float Precalcolato} & \textbf{Intero TF Compresso} \\
\midrule
\textbf{Occupazione Spazio} & $32$ o $64$ bit per posting (\texttt{float}/\texttt{double}). & Intero compresso (spesso $< 8$ bit medi per posting). \\
\addlinespace
\textbf{Comprimibilit\`a} & Scarsa o nulla: mantissa float a entropia quasi-casuale. & Elevatissima: concentrazione statistica su delta piccoli. \\
\addlinespace
\textbf{Flessibilit\`a Ranking} & Nulla: vincola l'indice alla specifica formula usata a build-time. & Totale: consente di calcolare qualsiasi funzione (BM25, varianti) a runtime. \\
\addlinespace
\textbf{Throughput CPU} & Nessun calcolo a runtime (lettura diretta). & Pochi cicli ALU per posting per calcolare la formula. \\
\bottomrule
\end{tabularx}
\caption{Confronto di sistema tra memorizzazione float e interi compressi.}
\label{tab:float-vs-int}
\end{table}

\begin{noteblock}{Risoluzione del Vincolo Memory-Bound}
Come analizzato nel \chapref{chap:efficienza-cpu}, i sistemi di Information Retrieval su larga scala sono rigorosamente \textbf{memory-bound}: il vero fattore limitante delle prestazioni \`e la banda di trasferimento dalla memoria principale (DRAM) ai registri della CPU ($\approx 80\text{ ns}$ di latenza), non il throughput dell'unit\`a aritmetica (ALU, $0.5\text{ ns}$). Ridurre la dimensione dei posting attraverso la compressione degli interi abbatte il volume di byte trasferiti sul bus di memoria; il modesto costo computazionale richiesto per valutare la formula a runtime (es. BM25) viene ampiamente compensato dal risparmio di banda.
\end{noteblock}

\FloatBarrier
\section{Tipologie di Query, Indici Posizionali e Riscrittura del Testo}
\label{sec:query-expansion-posizionali}

L'indice invertito essenziale a sole frequenze $(\text{DocID}, \text{tf})$ non \`e sufficiente per supportare ricerche complesse che richiedono la contiguit\`a delle parole o la riscrittura semantica delle interrogazioni. Questa sezione approfondisce la struttura dell'indice posizionale per l'elaborazione delle \textit{phrase queries} e illustra le tecniche di \textit{query expansion} e normalizzazione a tempo di esecuzione.

\FloatBarrier
\subsection{Phrase Queries e la Necessit\`a dell'Indice Posizionale}
\label{subsec:phrase-queries-indice-posizionale}
Nelle ricerche per frase esatta (\textit{Phrase Queries}), la query \`e racchiusa tra virgolette (es. \texttt{"neural network architecture"}). L'utente impone che i termini compaiano nel medesimo documento, nello stesso ordine e in posizioni adiacenti:
\begin{equation}
\text{pos}(t_{i+1}) = \text{pos}(t_i) + 1.
\end{equation}
L'indice invertito standard a sole frequenze $(\text{DocID}, \text{tf})$ non permette di validare questo vincolo, avendo scartato la posizione delle parole. Per supportare le phrase queries, la struttura viene estesa a \textbf{Indice Posizionale}, memorizzando per ogni posting la lista degli offset ordinali:
\begin{equation}
t \longrightarrow \big[ \langle \text{DocID}_k, \text{tf}_k, [\text{pos}_1, \text{pos}_2, \dots, \text{pos}_{\text{tf}}] \rangle \big].
\end{equation}

\begin{example}[Intersezione Posizionale per Phrase Match]
Si consideri la query \texttt{"neural network"}:
\begin{itemize}
    \item Posting per \texttt{"neural"}: $\dots \to \langle \text{Doc } 14, \text{tf}=2, \text{pos}=[4, 21] \rangle \to \dots$
    \item Posting per \texttt{"network"}: $\dots \to \langle \text{Doc } 14, \text{tf}=3, \text{pos}=[5, 89, 102] \rangle \to \dots$
\end{itemize}
L'algoritmo esegue l'intersezione su $\text{DocID} = 14$ e confronta le liste posizionali: poich\'e $\text{pos}(\texttt{neural}) = 4$ e $\text{pos}(\texttt{network}) = 5$, la condizione $\text{pos}(t_{i+1}) = \text{pos}(t_i) + 1$ \`e soddisfatta e il documento viene validato.
\end{example}
\textit{Overhead spaziale:} L'inclusione delle posizioni incrementa l'impronta complessiva su disco dell'indice invertito di un fattore compreso tra $2\times$ e $3\times$.

\FloatBarrier
\subsection{Moduli di Query Expansion e il Caso Studio ``Babbo Francesco''}
\label{subsec:query-expansion-babbo-francesco}
I motori di ricerca industriali intercettano la query grezza dell'utente attraverso un modulo di \textbf{Query Expansion} e riscrittura, con l'obiettivo di correggere refusi, risolvere acronimi ed espandere sinonimi contestuali. 

Tuttavia, espansioni lessicali non supervisionate basate su matrici di co-occorrenza o regole euristiche possono innescare gravi derive semantiche (\textit{semantic drift}), come illustrato dal celebre caso riscontrato nei motori commerciali in lingua italiana:
\begin{enumerate}
    \item \textbf{Query Utente:} L'utente inserisce la query a intento personale/familiare \texttt{"babbo francesco"}.
    \item \textbf{Espansione Sinonimica:} Il modulo associa il termine regionale \textit{babbo} al sinonimo \textit{pap\`a}.
    \item \textbf{Normalizzazione e Troncamento Accenti:} La pipeline rimuove gli accenti trasformando \textit{pap\`a} in \textit{papa} (omografo non accentato che designa il Pontefice).
    \item \textbf{Query Riscritta all'Indice:} La disgiunzione risultante inviata al motore diventa \texttt{("babbo" OR "pap\`a" OR "papa") AND "francesco"}.
    \item \textbf{Scoring \& Authority Ranking:} L'entit\`a globale \textit{"Papa Francesco"} possiede un'autorevolezza a priori (PageRank) e una co-occorrenza web sterminata rispetto a qualunque testo privato.
    \item \textbf{Risultato Top-1 Restituito:} La SERP collassa restituendo in prima posizione la voce enciclopedica su \textit{Papa Francesco (Jorge Mario Bergoglio)}, stravolgendo completamente l'intento informativo dell'utente.
\end{enumerate}

\FloatBarrier
\subsection{Stemming vs. Query Expansion}
\label{subsec:stemming-vs-expansion}
\`E fondamentale chiarire la linea di demarcazione tra normalizzazione morfologica ed espansione semantica:
\begin{itemize}
    \item \textbf{Stemming e Lemmatizzazione:} Operano riduzioni sintattiche all'interno della medesima radice etimologica (es. \textit{running}, \textit{runs} $\to$ \textit{run}, come descritto in Sezione~\ref{subsec:stemming-lemmatizzazione}). Non possono collegare vocaboli con radici distinte aventi identico significato (es. \textit{babbo} e \textit{pap\`a}).
    \item \textbf{Approccio Moderno:} I sistemi scalabili contemporanei riducono l'uso di stemming aggressivo in fase di indicizzazione (per preservare la specificit\`a semantica dei termini nell'indice) e gestiscono le varianti flessive e sinonimiche a runtime tramite modelli linguistici neurali di espansione contestuale.
\end{itemize}

\begin{noteblock}{Dall'Indicizzazione al Query Processing}
Costruite le strutture dati dell'indice e definite le modalit\`a di rappresentazione delle interrogazioni, il motore deve eseguire il recupero dei documenti rilevanti in tempo reale ($\le 50\text{--}100\text{ ms}$). I modelli teorici e i paradigmi algoritmici per l'attraversamento delle posting list, l'aggregazione dei punteggi, l'ottimizzazione tramite salti rapidi (\textit{Term-at-a-Time} e \textit{Document-at-a-Time}) e la loro realizzazione empirica in Python sono formalizzati e analizzati esaustivamente nel Capitolo~\ref{chap:indici-inversi-query-processing}.
\end{noteblock}
